\chapter*{Введение}
\addcontentsline{toc}{chapter}{Введение}
\label{ch:intro}

    Целью лабораторных работ №1.1 и №1.2 является изучение системы MATLAB для инженерных, финансовых и математических расчётов, а также освоение матричных операций и основ программирования функций.

    В лабораторной работе №1.1 рассматриваются: задание функции одной переменной, построение графиков функции, её производной и интеграла, графическое и аналитическое решение уравнения вида $a \cdot x + b = f(x)$, построение трёхмерного графика функции двух переменных $F(x,y)$.

    В лабораторной работе №1.2 рассматриваются: генерация векторов-столбцов и векторов-строк, перемножение векторов, вычисление среднего значения и дисперсии матрицы, программирование пользовательской функции в MATLAB.

    Все расчёты выполнены в соответствии с вариантом №24 (функция одной переменной) и вариантом №24 (размерности матрицы $I=30$, $J=26$).

\endinput